\scanpage{38}
% archive sheet number: 37
\mspage{( 3 )}
% L1–5
\emph{Application aux corps de nombres algébriques.} --- Soit $\mathrm A$ un corps de nombres algébriques; à toute valeur absolue $v$ de $\mathrm A$, on associe le complété $\mathrm A_{v}$ de $\mathrm A$ pour cette valeur absolue. D'où une application
\[
\mathrm Q(\mathrm A)\to\prod\mathrm Q(\mathrm A_{v})
\]
et cette application est injective (théorème de Hasse).

% L6–9
Si $v$ est une valuation discrète, toute classe de forme sur $\mathrm A_{v}$ est déterminée par $\varepsilon$, $w^{1}$ et $w^{2}$ $(^{4})$. Si $\mathrm A_{v}=\mathbf R$, le corps des réels, alors $\mathrm H(\mathbf R)$ s'identifie à $\mathbf Z/2\,\mathbf Z\,[t]$. Toute classe de forme est définie par deux nombres entiers $m$ et $n$ ($m$, nombre de carrés positifs; $n$, nombre de carrés négatifs).

% L10–11
Posant $n=\displaystyle\sum n_{j}\,2^{j}$ $(n_{j}=0\text{ ou }1)$, on trouve
\[
w^{j}(q)=\binom{n}{j}\bmod 2;\qquad
\text{en particulier}\quad w^{2^{j}}=n_{j}.
\]

% L12–15
Donc une classe de forme sur $\mathbf R$ est complètement déterminée par $w$. De tout ce qui précède on déduit le diagramme commutatif suivant :
\diagram{\begin{tikzcd}[ampersand replacement=\&]
\mathrm Q(\mathrm A)\arrow[r]\arrow[d,"w"] \& \displaystyle\prod\mathrm Q(\mathrm A_{v})\arrow[d,"{(w_{v})}"]\\
\mathrm H(\mathrm A)\arrow[r] \& \displaystyle\prod\mathrm H(\mathrm A_{v})
\end{tikzcd}}

% L16–18
On démontre ainsi que l'application $w$ est injective : Sur un corps de nombres toute classe de forme quadratique est déterminée par son rang et ses classes de Stiefel-Whitney.

% L19–22
\emph{Remarque.} --- On ne sait pas si l'application
\[
\operatorname{grad}^{2}(\mathrm Q(\mathrm A))\to\mathrm H^{2}(\mathrm A,\mathbf Z/2\,\mathbf Z)
\]
est toujours injective, et l'on ne sait pas non plus si tout élément d'ordre $2$ du groupe de Brauer de $\mathrm A$ est somme d'éléments de la forme $a.b$.

% L23–26
{\small
$(^{1})$ N. \textsc{Bourbaki}, \emph{Algèbre}, chap. IX, Hermann, Paris, 1959.\\
$(^{2})$ A. \textsc{Grothendieck}, \emph{Bull. Soc. math. Fr.}, 86, 1958.\\
$(^{3})$ T. A. \textsc{Springer}, \emph{Akad. Amst. Proc.}, 62 A, 1959.\\
$(^{4})$ E. \textsc{Witt}, \emph{Journal de Crelle}, 176, 1937.\par}
